%% notes-tex/multimodal/043-training-plan/sections/1-why.tex
%% SOURCE: results/overnight_2026_10_22_window_means.json (jobs 6136, 6137, 6150, 6152, seed 0,
%% registered window) as tabulated in notes-tex/multimodal/043-training-design/tables/t1-measured-rows.tex;
%% design review tables/t2-batch-share-rows.tex (expected weights); the 040 final-push list in
%% memory 040-final-push-plan-2026-10-22 (item 3, pooled-then-target); the manuscript framing from
%% paper/nature-biotech/sections/results.tex (Fig. 3 title) and sections/introduction.tex.
%% Owner quotes are verbatim from the 2026-10-23 direction relayed with this task.

\section{Why two stages\secstatus{todo}}
\label{sec:why}

\notesec{What one joint objective measured}
The overnight round of 2026-10-22 trained one objective over the nine datasets of
\texttt{joint\_9} and asked of it a per-target question: does the joint run beat a run on
Kemmeren 2014 alone at the same number of optimizer steps (\dsec{5.4})? At seed 0, on the
validation fold, the Kemmeren window mean read 0.078 for the records sampler against 0.110
alone, a transfer of $-0.032$; $-0.020$ for the pinball arm; and $+0.009$ for the balanced
sampler. On the 72 Sameith 2015 doubles the alone run read 0.335 and the three joint runs
0.208 to 0.230. One seed is below the decision rule's six, and no run of that round logged a
validation loss, so the comparison rests on window means and the training loss and cannot be
read for overfitting (\dsec{6}).

\notesec{Why the answer depended on the batch}
In a single objective the transfer to a target is decided by how much of each batch the
target holds. Under equal weighting the records sampler gives Kemmeren an expected weight of
0.066 and the promoter library 0.573; the balanced sampler gives every dataset 0.111
(design Table~2). The one positive transfer came from the sampler that moved Kemmeren's weight
up. Hypothesis (untested): the sign of the overnight transfer follows the target's expected
weight, not the value of the shared representation; one seed cannot separate the two.

\notesec{The same lesson in 040}
The environment case study (experiment 040) built its trainer around one target screen, the
others entering as auxiliary sources, and its final list of experiments adds a
pooled-then-target arm: train on every screen, then drop the auxiliary sources for the last 10
to 20\% of epochs. The direction drawn from both experiments is not to gear a trainer to one
target, and to ``pretrain'' over all phenotypes before fine-tuning for those of interest. For
043 the owner stated it on 2026-10-23: ``we want to just be able to train models over all
data. then we can eval them at the end, but then we can also fine tune on perturbation data to
see if it helps with sameith doubles for example.''

\notesec{Two questions, two stages}
The joint objective answers two questions at once. Whether one representation carries every
readout is a property of the pretrained model, read per dataset at the end of stage 1
(Sec.~\ref{sec:stage1}). Whether that representation helps a given target is a property of
the starting point, read by fine-tuning from it against training from scratch at matched
optimizer steps (Sec.~\ref{sec:stage2}). The stage-1 trainer then needs no target, and the
target-specific comparison moves to stage 2, where every batch belongs to the target.

\notesec{The paper's framing}
The manuscript's Fig.~3 states that one latent cell embedding generalizes across
systems-biology phenotypes (expression, morphology and fitness), and its introduction places
the work beside foundation cell models that pretrain on single-cell expression matrices and
fine-tune for phenotype, noting that industrial microbes lack such matrices but have diverse
perturbation data. Stage 1 is that claim's test on the 043 readouts, with morphology held out
of training and read as a probe; stage 2 is the fine-tuning half.

\subsection{Symbols added by this plan\secstatus{todo}}
\label{sec:symbols}

Every symbol of the design review (\dsec{1}, Table~1) keeps its meaning. The plan adds the
symbols of Table~\ref{tab:symbols}, alphabetical by letter name.

%% SOURCE: hand-authored notation contract; each row names the equation of this document that
%% defines it.
\begin{table}[htbp]
\centering
\footnotesize
\setlength{\tabcolsep}{4pt}
\caption[]{\textbf{Symbols added by the plan.} Symbols of the design review are not
repeated; ``Eq.'' is the equation of this document that defines the symbol.}
\label{tab:symbols}
\begin{tabular}{>{\raggedright\arraybackslash}p{34mm}>{\raggedright\arraybackslash}p{120mm}>{\raggedright\arraybackslash}p{16mm}}
\toprule
Symbol & Meaning & Defined in \\
\midrule
$\alpha$ & sampler exponent: dataset $s$ is drawn for a batch slot with probability $\pi_s(\alpha) \propto n_s^{\alpha}$; $\alpha = 0$ is the balanced sampler & Eq.~\ref{eq:sampler} \\
$\mathrm{bg}(b),\ \mathcal{P}^{\mathrm{bg}}_{\mathrm{bg}(b)}$ & background strain of record $b$ (S288C, BY4741, BY4742, FY4, FY5, Y8205); the fixed marks that strain carries relative to S288C & Eq.~\ref{eq:bgmarks} \\
$\Delta^{\mathrm{ft}}_{\mathrm{tot}},\ \Delta^{\mathrm{ft}}_{k}$ & fine-tune gain against training from scratch for $k_1 + k$ steps, and for $k$ steps & Eq.~\ref{eq:ftgain} \\
$\mathbf{E}^{\mathrm{nbr}}_{k},\ \mathbf{E}^{\mathrm{orig}}$ & learned vector of an empty relation-$k$ neighborhood; learned table of mark origin (engineered, background) & Eqs.~\ref{eq:gather},~\ref{eq:bgmarks} \\
$\bar h^{(k)}_{b,i}$ & mean perturbed state of the relation-$k$ neighbors of gene $i$ in record $b$ (readout gather) & Eq.~\ref{eq:gather} \\
$\iota_t$ & origin of mark $t$: 0 engineered, 1 background & Eq.~\ref{eq:bgmarks} \\
$k_1,\ k,\ k_{\mathrm{tot}}$ & optimizer steps of stage 1, of stage 2, and their sum & Eq.~\ref{eq:ftruns} \\
$\mathrm{lr}_2,\ \mathcal{F}$ & stage-2 learning rate; freezing policy (none, trunk frozen, heads only) & Sec.~\ref{sec:ftstart} \\
$\pi_s(\alpha)$ & probability that a batch slot draws dataset $s$ & Eq.~\ref{eq:sampler} \\
$\rho_b$ & agreement of the route-A and route-B predicted profiles of isolate $b$, centered & Eq.~\ref{eq:routecons} \\
$\theta^{(1)}$ & all trainable parameters at the last epoch of a stage-1 run (the window-end checkpoint) & Eq.~\ref{eq:ftruns} \\
$W_{\mathrm{ga}}$ & projection of the nine gathered neighbor states, $\in \mathbb{R}^{d \times 9d}$ & Eq.~\ref{eq:gather} \\
\bottomrule
\end{tabular}
\end{table}
